\chapter{仮想環境のセットアップ}

\section{フォルダに移動する}

まず課題フォルダに移動します。

\begin{wslbox}[課題フォルダへ移動]
\begin{lstlisting}[style=bash]
$ cd \
  ~/NMR-lab-kadai/2026_B4_assignment
\end{lstlisting}
\end{wslbox}

移動できたか確認します。

\begin{wslbox}[現在地の確認]
\begin{lstlisting}[style=bash]
$ pwd
/home/username/NMR-lab-kadai/2026_B4_assignment
\end{lstlisting}
\end{wslbox}

このように表示されれば OK です。

\section{uv sync を実行する}

\begin{wslbox}[仮想環境のセットアップ]
\begin{lstlisting}[style=bash]
$ uv sync
\end{lstlisting}
\end{wslbox}

初回実行時は Python のダウンロードと全ライブラリのインストールが
自動で行われます。回線速度によっては数分かかる場合があります。

\subsection{実行中の出力例}

\begin{wslbox}[uv sync の出力例]
\begin{lstlisting}[style=bash]
Using CPython 3.11.x
Creating virtual environment at: .venv
Resolved 47 packages in 0.5ms
Installed 47 packages in 8.2s
 + bm3d==4.0.2
 + ipywidgets==8.1.5
 + japanize-matplotlib==1.1.3
 + jupyterlab==4.3.6
 + matplotlib==3.9.4
 + numpy==1.26.4
 + PyWavelets==1.5.0
 + scikit-image==0.22.0
 + scipy==1.13.1
 ...
\end{lstlisting}
\end{wslbox}

このような出力が出れば完了です。

\subsection{2回目以降の実行}

すでに環境が揃っている場合は、\cmd{uv sync} を再実行しても
差分だけが更新されます。不足しているパッケージがあれば追加されます。

\begin{wslbox}[2回目以降の出力例]
\begin{lstlisting}[style=bash]
$ uv sync
Resolved 47 packages in 0.5ms
Audited 47 packages in 1ms
\end{lstlisting}
\end{wslbox}

\section{セットアップの確認}

正しくインストールされたか確認します。

\begin{wslbox}[インストール済みパッケージの確認]
\begin{lstlisting}[style=bash]
$ uv pip list
Package                  Version
------------------------ -------
bm3d                     4.0.2
japanize-matplotlib      1.1.3
jupyterlab               4.3.6
matplotlib               3.9.4
numpy                    1.26.4
...
\end{lstlisting}
\end{wslbox}

\cmd{numpy}、\cmd{jupyterlab}、\cmd{bm3d} などが表示されれば
正しくインストールされています。

\begin{tipbox}[uv.lock の役割]
\cmd{uv.lock} ファイルがあることで、
先輩・後輩・教員の全員がまったく同じバージョンのライブラリを使えます。
「自分の環境では動いたのに先輩の環境では動かなかった」という問題を防ぐための仕組みです。
\end{tipbox}
